%
% This is a template for a new research paper. For information on this file
% please contact Joe Loughry at Tel. +1 720 277 7800 (time zone GMT minus 7
% hours) or Email: Joe.Loughry@cs.du.edu or mailto:Joe.Loughry@gmail.com
%

\documentclass[conference]{IEEEtran}

\usepackage{cite}

\usepackage[english,british]{babel}
\usepackage{graphicx}
\usepackage[binary-units]{siunitx}
\usepackage[obeyspaces,hyphens]{url}
\newcommand{\URL}[1]{$\langle$\url{#1}$\rangle$}
\usepackage[plainpages=false,pdfpagelabels]{hyperref}

\hyphenation{op-tical net-works semi-conduc-tor}

\begin{document}

\title{Optical TEMPEST
}

\author{\IEEEauthorblockN{Joe Loughry}
\IEEEauthorblockA{University of Denver, Department of Computer Science \\
Denver, Colorado 80208 USA \\
Email: joe.loughry@cs.du.edu}}

\maketitle

\begin{abstract}
	Research on optical TEMPEST has exploded since 2002 when the first papers
appeared from widely separated locations within a week of each other. In the
intervening time, vulnerabilities have evolved along with hardware, and
several new threat vectors have appeared. Although the supply chain ecosystem
of Ethernet has reduced the vulnerability of billions of devices through use
of standardised PHY chips, other recent trends including High Frequency
Trading (HFT) in the financial sector, the Internet of Things (IoT), and
inexpensive drones have made it relevant again.

\end{abstract}

\section{Introduction}

Since the publication sixteen years ago of the first two papers on optical
TEMPEST, open source information on security of compromising emanations---and
side channels in general---has exploded. Before the late 1990s, only a
handful of papers had been published in the open literature on the subject of
TEMPEST\footnote{TEMPEST, here, is the U.S.\ National Security Agency (NSA)
code name for `the problem compromising radiation'---radio frequency (RF)
\emph{or} acoustic, according to the original reference, which was only
declassified in 2007---including both exploitation and control
\cite{NSATempest2007}. Beyond the 1972 definition of TEMPEST, military and
academic researchers have expanded the spectrum to include DC, optical,
thermal, magnetic, and acceleration side channels.} and the topic was mostly
relegated to the folklore of infosec; the only two scientific studies of
unintentional compromising emanations remained van Eck \cite{vanEck1985} and
Smulders \cite{Smulders1990} although Wright's book around the same time
\cite{Wright1987} described anecdotal reports dating back to the first world
war. But following Kocher's seminal 1996 paper on side channel attacks
\cite{Kocher1996} and Kuhn \& Anderson \cite{Kuhn1998a} showing that covert
channels were not limited by the system boundary, two papers appeared within
a week of each other on complementary aspects of optical emanations
\cite{Kuhn2002,Loughry2002a}. Since then, the two papers, between them, have
been cited more than 300 times.

\section{Stories never told before}

Optical TEMPEST was discovered at a bank. In the raised-floor computer room
on the sixth level of a glass high-rise building in downtown Seattle,
surrounded by other glass high-rises with their own computer rooms visible at
night by the reddish glow of LED indicators, I was working late. Dial-up
modems had not yet gone completely extinct, \SI{10}{\mega\bit\per\second}
Ethernet was increasingly common on PCs, and leased lines ran everywhere from
the computer room to branch offices, thousands of them. This was the
environment where optical TEMPEST was discovered in 1992. Cautious
experiments were performed, a literature survey was quietly done to see if
anyone had ever noticed it before, and the National Computer Security Center
(NCSC) was asked if they knew of it. All inquiries ran into a classified
information roadblock; just about all that was known at the time, in the open
literature, was that nearly everything about TEMPEST was classified (but it
was probably not an acronym).

\subsection{Missed opportunities}

Around this time, I made a serious mistake. Unbeknown to me, Markus Kuhn, a
postgraduate student at Cambridge, was working along similar lines. I had
thought of optical emanations from video display screens but dismissed the
idea as physically impossible---without ever testing it. I thought the decay
time of Cathode Ray Tube (CRT) phosphors was too slow to carry information
about the video signal in the diffuse light available to non-line-of-sight
interception, and consequently, I never looked.

Markus Kuhn found and successfully exploited a tiny ripple near the peak of
the response curve of video display phosphors. The gross shape of the curve
belies the fact that that tiny ripple is detectable in the time domain
optical signal, if your detector is fast enough.

Kuhn's detector---a photomultiplier tube---was better than mine; the
gain--bandwidth product is better than that of the large-area
photodiode/transimpedance amplifier combination I used, but we were looking
for different things: Kuhn for diffuse emanations from an entire screen, and
I for line-of-sight photons from a particular LED (although we did eventually
figure out how to separate multiple superimposed signals from diffuse optical
emanations collected by non-line-of-sight means). The difference is that
Markus Kuhn actually looked, and he found an effect I
missed.\footnote{Speaking of video signals, laser printers intensity-modulate
an infrared laser with a video signal, and some plastics are transparent to
infrared light. It might be worthwhile to measure laser printers for
information-bearing optical emanations outside visible wavelengths, as Kubiak
has done for RF \cite{Kubiak2014,Kubiak2017,Kubiak2017b,Kubiak2017c},
Ula\c{s} {\it et al.}\ have done for conducted powerline emanations
\cite{Ulas2016}, and Enev {\it et al.}\ have done for conducted powerline
emanations from video displays \cite{Enev2011}.}

\subsection{Delay of first publication}

Why didn't we publish in 1994? Part of the reason was perfectionism; I wanted
to be able to explain the phenomenon and make predictions, not only describe
it. In addition, by that time I had left the bank and was working for a
defence contractor on a classified project. I had a security clearance now,
and a greatly expanded awareness of counterintelligence sources and methods.

And so, following procedures, we submitted the paper to NSA for approval to
publish. It took a year and a half to approve.\footnote{Actually, it didn't
quite happen that way. Our paper was approved for publication very quickly,
in only a few weeks; we submitted it to the 10th USENIX Security Symposium
where it was immediately accepted. A few days later, NSA called us back, and
in a panic, insisted we withdraw the paper from the conference. I had to
apologise to the programme chair; it was terribly embarrassing, and the delay
in publishing was two years.} Eventually, NSA wrote back and said, `approved
for public release'. I wonder what they spent all that time doing.

\subsection{Roads not taken}

Universal Serial Bus (USB) devices hadn't happened yet, but today they're
ubiquitous. Later and informal investigation of some light-up USB cables---a
fad that didn't last---turned up evidence of Class II optical emanations
related to data passing through the USB cable, but evidence for Class III
optical emanations was inconclusive.

Considered narrowly outside the scope of optical TEMPEST, for the purpose of
this paper, are line-of-sight attacks that essentially reduce to direct or
indirect imaging of the display
\cite{Backes2008,Balzarotti2008,Backes2009a,Raguram2011,Xu2013a,Jenkins2013a}.
This is somewhat unfair, as the original attack by Loughry \& Umphress
required, for the most part, line-of-sight---except, as previously mentioned,
in \S 8.2---but in the time domain, not space. Optical TEMPEST is a time
domain effect. Other remote attacks employing optical means, such as visual
or interferometric measurement of keyboard or printer acoustic emanations
\cite{Asonov2004,Zhuang2005,Berger2006,Backes2010}, or the highly novel
reverse covert channel using a scanner \cite{Nassi2017a}, are outside the
scope of this paper. Finally, induced optical emanations
\cite{Sepetnitsky2014a,Guri2016b,Guri2017a,Guri2017b,Lopes2017a,Guri2017c,
Zhou2017,Zhou2018a} and \cite[Appendix A]{Loughry2002a} are properly
considered out-of-band covert channels
\cite{Lampson1973,Hanspach2014,Carrara2016}, despite being compromising
optical emanations in the time domain, because they are purposely induced by
a nefarious software or hardware agent, or by activity controllable by a
third party, introduced into the target system by the attacker.

\section{MAC and PHY}

We looked hard, at the time, at the other two large populations of blinking
LEDs in the world: disk activity lights and Network Interface Cards (NICs).
Neither source proved fruitful; we were unable to find any evidence of Class
III optical emanations from storage devices or link activity indicators on
Ethernet cards. The one exception---and it was a bad one---was WAN interfaces
on the back panel of enterprise routers, devices which live in racks that
sometimes back up to windows; see \cite[\S 4.3.1]{Loughry2002a}. Aside from
those, the complete absence of compromising optical emanations from Ethernet
link activity LEDs is believed to be a consequence of the fact that the
Ethernet protocol is well divided into two layers: MAC and PHY.

In the Ethernet protocol, the Media Access Control layer (MAC) marshals bits
into frames and hands them off to the Physical layer (PHY), which deals
exclusively with voltages and waveforms and wires, or radio, or fibre optics.
In the case of twisted pair wire, only a few suppliers make PHY chips and
generally they do it right, providing dedicated pins on the PHY chip for
connecting LEDs for status indication and internally stretching pulses to the
LEDs to make high speed activity visible to humans. The contrast with the
situation we found regarding relatively low-speed serial interfaces is stark;
there, the temptation was overwhelming for circuit designers to drive LEDs
directly from generously high voltage and high current serial communication
signals---arguably providing reliable indication of signal quality and
marginal signal levels at very low cost.

But see \S forward reference to High Frequency Trading (HFT)

bespoke hardware for ultra-low-latency Ethernet

HFT

Optical TEMPEST began in a bank, developed in counterintelligence,
and now circles back to fintech. FPGA for latency \cite{Hurd2018a}.

The first part of this paper is a survey of further work done by others since
the publication of \cite{Loughry2002a} in 2002. This is followed by a
description of the design process for a new product according to these
principles.

\section{Design of a new product with optical TEMPEST principles
in mind}

According to these hard-won principles.

\subsection{Simple photoconductive photodiode design}

The photodiode circuit is shown in Figure \ref{figure:photodiode_pullup}.
The photodiode operates in reverse bias (photoconductive) mode for two
reasons: speed, and simplicity of implementation. The same photodiode
operated in photovoltaic mode would be more sensitive to very low level
signals and have a lower dark current, but would require a transimpedance
amplifier for current-to-voltage conversion, which would make the design
more complicated and thereby more difficult to evaluate for security.

% Use [!t] for figures in IEEEtrans papers..
\begin{figure}[!t]
    \centering
	\includegraphics[height=2in]{graphics/photodiode_pullup_and_GPIO_protection.png}
	\caption{The photodiode circuit is purposely made as simple as possible
        for transparency of implementation; the \SI{10}{\kilo\ohm} pull-up
        resistor for reliability and a \SI{100}{\ohm} series resistor to
        protect the general-purpose IO (GPIO) pin from being shorted to
        ground in case it were accidentally set to output a HIGH logic level
        at the same time the photodiode was illuminated.}
	\label{figure:photodiode_pullup}
\end{figure}

% See bare_conf.tex for example of a two-column subfigure in IEEEtrans, and
% for the preferred table format.

Risk of acoustic information leakage from acousto-optic modulator or MMD
modulator. Cooling fan speed modulation acoustic information leakage risk.
Light-tight design. RF, acoustic, vibration, temperature, ELF, power line
conducted emissions: see lit.\ survey in recent papers for a comprehensive
list.

To be honest, I have reservations about the lack of galvanic isolation
through the metal of this component.

\section{Energy Gapping}

Cite the `energy gap' principles from Clive Robinson.

Conductive tape on the seams and for isolation between high side and low side
inside.

\section{Conclusion}

Since the publication in 2002 of the first peer-reviewed research on
compromising optical emanations, other researchers have carried it further.
But technological progress has shifted the boundaries of what was possible,
necessitating re-visit of the same vulnerabilities from time to time. This is
a general principle of security; vulnerabilities, once fixed, sometimes do
not stay fixed.

\section*{Acknowledgments}

Thanks to many anonymous and pseudonymous commenters on Hacker News
\URL{news.ycombinator.com} for helping me develop the ideas in this paper.

% trigger a \newpage just before the given reference
% number - used to balance the columns on the last page
% adjust value as needed - may need to be readjusted if
% the document is modified later
%\IEEEtriggeratref{8}
% The "triggered" command can be changed if desired:
%\IEEEtriggercmd{\enlargethispage{-5in}}

\bibliographystyle{IEEEtran}
% argument is your BibTeX string definitions and bibliography database(s)
\bibliography{IEEEabrv,consolidated_bibtex_file}

\end{document}

